\documentclass[12pt]{article}


% --- math/fonts ---
\usepackage{amsmath,amssymb}
\usepackage{braket}
\usepackage[T1]{fontenc}
\usepackage{libertine}
\usepackage[utf8]{inputenc}

% --- graphics/tikz ---
\usepackage{tikz}
\usepackage{qcircuit}

% --- other stuff ---
\usepackage[margin=1in]{geometry}
\usepackage{pdfpages}
\usepackage{multicol}
\usepackage[shortlabels]{enumitem}
\usepackage{cancel}
\usepackage{graphicx}
\usepackage{float}
\usepackage{listings}
\usepackage{booktabs}
\lstset{
    language=Python,
    basicstyle=\ttfamily\footnotesize,
    keywordstyle=\bfseries,
    frame=single,
    breaklines=true,
    keepspaces=true,
    columns=fullflexible,
    showstringspaces=false
}
\let\epsilon\varepsilon
\usepackage{enumitem}
\setlist[enumerate]{label=\arabic*.}
\usepackage[authoryear,round]{natbib}
\usepackage[hidelinks]{hyperref}


\newcommand{\class}{CSCE883 Machine Learning}           
\newcommand{\examnum}{Project Proposal}      
\newcommand{\examdate}{10/18/2026} % Fill in the proposal date.        
\newcommand{\timelimit}{}             




\begin{document} 
\pagestyle{plain}
\thispagestyle{empty}

\noindent
\begin{tabular*}{\textwidth}{l @{\extracolsep{\fill}} r @{\extracolsep{6pt}} l}
 \textbf{\class} & \textbf{Name:} & \textit{Taylor Belcher}\\       
\textbf{\examnum} &&\\
\textbf{\examdate} &&\\
\end{tabular*}\\
\rule[2ex]{\textwidth}{2pt}
% ---

% Write your project proposal here.

My research area is in formal verification methods with functional programming languages and quantum computing. In particular, I am interested in formalizing proofs from mathematics. The explosion of LLMs has also led to advances in capability in this area as models have medaled in the International Mathematics Olympiad. OpenAI recently announced a solution to a Millenium Prize problem (the Navier-Stokes conejctur) and also caused a stir in the mathematical community with the release of hundreds of proofs on many topics. These natural language proofs were released along with a formalization of the results, written the interactive theorem proving language Lean. I have been working with a language called Agda, but they are related. Formalization of natural language proofs predates the recent explosion in AI capability, but they are mutually helpful. LLM's ability to generate code means we can more quickly produce formalized mathematics, and the strict type requirements of formalization means we have good ways of building objective functions for reinforcement learning for AI models. 

Proof formalization as a field of computer science comes from a desire in the field of mathematics to make sure that our mathematics is correct. Even before the recent release of hundreds of results from OpenAI, understanding all of published mathematics and being confident that it was all correct was already a well-known problem. Formalized proofs in Lean (or another language) are meant to be a solution to this issue. Instead of trusting that other mathematicians and their peer-reviewers to do correct work, you only have to trust the kernel of the Lean programming language to not allow a proof of true=false. Although you are ultimately just moving where you put your trust, there is good reason to trust the kernel of Lean or another language like the one I work in, Agda. 

However, there is a meta-verification problem. How can you be sure that the formalization of a theorem and its proof in Lean is a faithful translation of the natural language proof? Now you not only have to read the natural language proof, but you also digest and verify the Lean code, which could be even longer and much more difficult to read since it is optimized for verification. (Although interactive theorem provers have long been working to improve on the balancing act between readability and formal correctness, like other normal programming languages.) A paper published in response to Open-AI's Navier-Stokes solution describes this problem exactly and gives examples from OpenAI's paper that identify where the Lean proof does not actually correspond to the natural language proof. It may be that the natural language proof is correct, but for the quantity of proofs and their length, we are not as justified in simply trusting the Lean certificate as we might hope. 

My proposal is to build a model that is rewarded for checking correspondence between formalized mathematics and natural language mathematics rather than generating either, although I am not sure if it will end up having to be able to generate expected translations in order to check faithful translations. My idea is that while we can reasonably trust the kernel of a language like Lean to faithfully judge whether a formalized proof is correct, we will need a separate judge to give us a confidence level on how well the formalized proof corresponds to the natural language proof. In terms of "is the theorem correct", this does not matter as much because if we trust the formal proof then we have evidence the theorem is correct, but we are also concerned about process and understanding in mathematics and we want to know what kinds of \emph{arguments} are useful, as we have seen that understanding arguments for one proof can be fruitful for other proofs. (Whether this becomes moot as mathematics stops being a human endeavor, I am unsure. But for now I am assuming that humans will still want to pursue frontier mathematics and understand, even if AI is better at it, much like people did not stop playing chess after Stockfish.)

Lean has become the most popular language for formalizing mathematics and there are data sets with pairs of natural language and Lean code available. I will start with Lean, but I am also interested in building something for Agda. Agda is far less popular and there are not nearly as many resources for data to train on. I may have to generate some data with LLMs and hand-label, but that will result in small datasets. Perhaps it would be possible to train on Lean and then use that model to specialize to Agda. 

\end{document}
